
Current practice in large language model trustworthiness evaluation treats reliability, robustness, and fairness as independent metrics. Benchmarks such as TruthfulQA, BOLD, and AdvGLUE measure these dimensions separately, reporting per-dimension scores without analyzing inter-dimensional relationships. This study tests whether trustworthiness dimensions correlate when measured synchronously on the same model outputs.

The evaluation landscape has developed frameworks that assess multiple dimensions together---TrustVis evaluates safety and robustness jointly, MLLMGuard reports scores across five dimensions, and BOLD benchmarks demographic bias---but these frameworks do not systematically compute or test cross-dimensional correlations. The present work fills this gap by treating evaluation logs as multi-dimensional datasets and analyzing correlation structure.

Synchronized evaluation operationalizes three constraints: (1) same model checkpoint across all measurements, (2) same input prompts, and (3) same generation parameters. This contrasts with sequential evaluation approaches that modify inputs between dimensions or benchmarks that evaluate dimensions on different datasets. Under synchronized evaluation, observed correlations reflect dimensional relationships rather than methodological artifacts.

On 817 TruthfulQA prompts with Llama-2-7b-chat, two correlation patterns were validated. Reliability and robustness correlate positively on factual prompts (r=0.72, 95\% CI [0.67, 0.77], p<0.001, n=343), interpreted as consistent with shared memorization: models that retrieve correct factual answers also retrieve them consistently across paraphrases. Fairness and reliability correlate negatively overall (r=-0.25, 95\% CI [-0.31, -0.18], p<0.001, n=817), interpreted as consistent with RLHF fine-tuning prioritizing demographic fairness at a cost to factual accuracy. A third hypothesis testing whether correlation magnitudes differ between factual and misinformation prompts was inconclusive (Fisher z-test p=0.788, n=10 per stratum), due to insufficient sample size.

The contributions are methodological and empirical. Methodologically, synchronized evaluation enables discovery of correlation structure in existing benchmarks without constructing new data. Empirically, two coupling patterns are validated with mechanistic interpretations linking correlations to training processes.
